\documentclass[../distribution_theory_notes.tex]{subfiles}
\begin{document}
\section{Aula 28 - 02 de Dezembro, 2024}

Previamente, finalizamos nossa apresentação da Transformada de Fourier, o que foi feito em três níveis:
\begin{itemize}
	\item[(1$^\circ$)] funções integráveis \(\bigl(L^1:=L^1(\mathbb{R}^n)\bigr)\);
	\item[(2$^\circ$)] funções de Schwartz \(\bigl(\mathcal{S}=\mathcal{S}(\mathbb{R}^n)\bigr)\); e
	\item[(3$^\circ$)] distribuições temperadas \(\bigl(\mathcal{S}':=[\mathcal{S}]'=\text{dual de }\mathcal{S}\bigr)\).
\end{itemize}

Quando \(f\in L^1\), sua transformada é a função \(\hat f\in \mathcal{C}_0\) dada por
\[
	\hat f(\xi):=\int_{\mathbb{R}^n}e^{-2\pi i x\cdot \xi}f(x)\,\mathrm{d}x,
	\qquad \xi\in\mathbb{R}^n,
\]
enquanto que sua transformada inversa é a \(\check f\in\mathcal{C}_0\) definida pondo
\[
	\check f(x):=\hat f(-x)=\int_{\mathbb{R}^n}e^{2\pi i x\cdot \xi}f(\xi)\,\mathrm{d}\xi,
	\qquad x\in\mathbb{R}^n.
\]
Note que \(\check f=R\circ\mathcal{F}(f)\), \(Rf(x)=f(-x)=\check f(x)\).

Demonstramos que:
\begin{itemize}
	\item[(i)] \(\mathcal{F}:\mathcal{S}\longrightarrow\mathcal{S}\) é linear contínuo \((f\in\mathcal{S}\Rightarrow\hat f\in\mathcal{S})\), de modo que \(\mathcal{S}\) é invariante pela transformada, o que pressupôs mostrar que \(\mathcal{S}\hookrightarrow L^1\);
	\item[(ii)] como consequência disso, e da igualdade
	      \[
		      \int \hat\varphi\,\psi=\int \varphi\,\hat\psi,
		      \qquad \varphi,\psi\in L^1,
	      \]
	      foi imediato estender \(\mathcal{F}\) a \(\mathcal{S}'\), pondo, para \(u\in\mathcal{S}'\) e \(\psi\in\mathcal{S}\),
	      \[
		      \langle \hat u,\psi\rangle:=\langle u,\hat\psi\rangle
		      \qquad
		      \bigl(\text{ou }\langle \overline{\mathcal{F}}(u),\psi\rangle
		      =\langle u,\mathcal{F}(\psi)\rangle\bigr).
	      \]
\end{itemize}

Neste ponto, destacamos que:
\begin{itemize}
	\item[(a)] \(L^p\hookrightarrow\mathcal{S}'\), \(1\leq p\leq\infty\), isto é, \(f\in L^p\Rightarrow f\) transformável;
	\item[(b)] \(\mathcal{S}'\hookrightarrow\mathcal{D}'\) (pois \(\mathcal{C}_c^\infty\hookrightarrow\mathcal{S}\) densamente).
\end{itemize}

\begin{itemize}
	\item[(iii)] \(\mathcal{F}:\mathcal{S}\longrightarrow\mathcal{S}\) é um isomorfismo com inversa
	      \[
		      \mathcal{F}^{-1}=R\circ\mathcal{F},
		      \qquad Rf(x)=f(-x),
	      \]
	      sendo \(R\) obviamente um isomorfismo de \(\mathcal{S}\).
\end{itemize}

A prova que demos baseou-se nas ``soluções da EDP do calor'', usando o ``ponto fixo'' da transformada
\[
	f_a(x)=e^{-\pi a|x|^2},\qquad a>0,
\]
para o qual
\[
	\hat f_a(\xi)=a^{-n/2}e^{-\pi|\xi|^2/a}.
\]
Com \(a=4\pi t\), nos dando
\[
	\hat f_{4\pi t}(\xi)
	=\frac{e^{-\frac{|\xi|^2}{4t}}}{(4\pi t)^{n/2}}
	=(2\sqrt{t})^{-n}\frac{e^{-\left|\frac{\xi}{2\sqrt{t}}\right|^2}}{\pi^{n/2}}
	=g_s(\xi),
\]
sendo
\[
	g(y)=\frac{e^{-|y|^2}}{\pi^{n/2}}
	\qquad\text{e}\qquad
	g_s(\xi)=s^{-n}g(\xi/s),\quad s=2\sqrt{t}.
\]
Dessa maneira, as soluções
\[
	u(x,t)=\bigl(\hat f_{4\pi t}*f\bigr)(x),
	\qquad x\in\mathbb{R}^n,\ t>0,
\]
e \(f\in\mathcal{S}\), de
\[
	u_t=\Delta u,
\]
cumprem
\[
	\lim_{t\to0^+}u(x,t)=f(x)
\]
nas funções contínuas e em \(L^1\).

Nosso estudo dessa extraordinária ferramenta deverá ser concluído hoje. O objetivo é estabelecer três fatos:
\begin{itemize}
	\item[(i)] o \underline{Teorema de Plancherel}: \(\overline{\mathcal{F}}:L^2\longrightarrow L^2\) está bem definido e é um ``isomorfismo unitário'' (exibindo mais um espaço invariante pela transformada);
	\item[(ii)] a \underline{transformada de Fourier--Laplace} de uma \(u\in\mathcal{E}'\), como preparatório para o terceiro;
	\item[(iii)] os \underline{Teoremas de Paley--Wiener}, os quais caracterizam as transformadas de
	      \begin{itemize}
		      \item[(1$^\circ$)] funções em \(\mathcal{C}_c^\infty\); e
		      \item[(2$^\circ$)] distribuições em \(\mathcal{E}'\).
	      \end{itemize}
\end{itemize}

\begin{tcolorbox}[
		skin=enhanced,
		title=Observação,
		fonttitle=\bfseries,
		colframe=black,
		colbacktitle=cyan!75!white,
		colback=cyan!15,
		colbacklower=black,
		coltitle=black,
		drop fuzzy shadow
	]
	Para o que segue, usaremos, provisoriamente, a notação \(\overline{\mathcal{F}}\) para a extensão de \(\mathcal{F}:\mathcal{S}\to\mathcal{S}\) a \(\mathcal{S}'\). Ao final da análise, abandonaremos esse cuidado.
\end{tcolorbox}

Mas, antes disso, vale a pena destacar:

\begin{crl*}
	\(\overline{\mathcal{F}}:\mathcal{S}'\longrightarrow\mathcal{S}'\) é um isomorfismo, cuja inversa é \(R\circ\overline{\mathcal{F}}\), \(R\psi:=\breve\psi\), isto é, \(\breve\psi(x)=\psi(-x)\), a reflexão.
\end{crl*}
\begin{proof*}
	Com efeito, a continuidade de \(R\circ\overline{\mathcal{F}}\) é clara, como composta de aplicações (lineares) e contínuas.

	Por outro lado, dados \(u\in\mathcal{S}'\) e \(\psi\in\mathcal{S}\), do teorema anterior,
	\[
		(\hat\psi)^\vee=\psi,
	\]
	donde
	\begin{align*}
		\bigl\langle (R\circ\overline{\mathcal{F}})\circ\overline{\mathcal{F}}(u),\psi\bigr\rangle
		 & =\bigl\langle \overline{\mathcal{F}}(\overline{\mathcal{F}}u),R\psi\bigr\rangle \\
		 & =\bigl\langle \overline{\mathcal{F}}(u),(R\psi)^\wedge\bigr\rangle              \\
		 & =\bigl\langle \overline{\mathcal{F}}(u),\breve{\hat\psi}\bigr\rangle            \\
		 & =\bigl\langle u,(\breve{\hat\psi})^\wedge\bigr\rangle
		=\langle u,\psi\rangle,
	\end{align*}
	pois \((R\psi)^\wedge=\breve{\hat\psi}\), conforme as propriedades de \(\mathcal{F}\) composta com rotações. De fato, se \(r(x)=-x\), então \(R\psi=\psi\circ r\), \((r^*)^{-1}=r\), e
	\[
		\breve{\hat f}=\hat f\circ r
		\Longrightarrow
		(\breve{\hat f})^\wedge
		=(\hat f\circ r)^\wedge
		=(\hat f)^\wedge\circ(r^*)^{-1}
		=(\hat f)^\wedge\circ r
		=(\hat f)^\vee=f. \text{ \qedsymbol}
	\]
\end{proof*}

\subsection{A transformada de Fourier em \(L^2\)}

Queremos determinar mais um espaço que seja invariante pela transformada. Os dois anteriores eram TVS's não normáveis e queremos agora apresentar um não só normável como também Hilbert, assim como ocorre com as séries de Fourier
\[
	\mathcal{F}:L^2(-\ell,\ell)\longrightarrow \ell^2(\mathbb{Z}),
\]
um isomorfismo unitário.

Para essa tarefa, vamos demonstrar ``gradualmente'':

\begin{lemma*}[Identidade de Parseval]
	Se \(f\in L^2\), então \(\overline{\mathcal{F}}f\in L^2\) e, além disso, vale a identidade
	\[
		\|\overline{\mathcal{F}}f\|_{L^2}=\|f\|_{L^2}.
	\]
	Mais ainda, se \(f,g\in L^2\), então
	\[
		(\overline{\mathcal{F}}f,\overline{\mathcal{F}}g)_{L^2}=(f,g)_{L^2},
	\]
	em que
	\[
		(f,g)_{L^2}=\int_{\mathbb{R}^n}f\,\overline g
	\]
	é o produto interno de \(L^2\).
\end{lemma*}
\begin{proof*}
	Primeiramente, dados \(\varphi,\psi\in\mathcal{S}\), provamos que
	\[
		(\#)\qquad (\hat\varphi,\hat\psi)_{L^2}=(\varphi,\psi)_{L^2}.
	\]

	\begin{tcolorbox}[
			skin=enhanced,
			title=Observação,
			fonttitle=\bfseries,
			colframe=black,
			colbacktitle=cyan!75!white,
			colback=cyan!15,
			colbacklower=black,
			coltitle=black,
			drop fuzzy shadow
		]
		Isso significa que
		\[
			\mathcal{F}:(\mathcal{S},\|\cdot\|_{L^2})\longrightarrow(\mathcal{S},\|\cdot\|_{L^2})
		\]
		é um operador unitário.
	\end{tcolorbox}

	Com efeito, sendo \(\mathcal{S}\hookrightarrow L^1\) e \(\mathcal{F}:\mathcal{S}\to\mathcal{S}\), podemos escrever
	\[
		(\hat\varphi,\hat\psi)_{L^2}
		=\int \hat\varphi\,\overline{\hat\psi}
		=\int \varphi\,[\overline{\hat\psi}]^\wedge.\tag{*}
	\]
	Como
	\begin{align*}
		\overline{\hat\psi(\xi)}
		 & =\overline{\left(\int e^{-2\pi i x\cdot\xi}\psi(x)\,\mathrm{d}x\right)}
		=\int \overline{e^{-2\pi i x\cdot\xi}\psi(x)}\,\mathrm{d}x                 \\
		 & =\int e^{2\pi i x\cdot\xi}\overline{\psi(x)}\,\mathrm{d}x
		=(\overline\psi)^\vee(\xi),
	\end{align*}
	porque a conjugação, como função \(\mathbb{C}\to\mathbb{C}\), é contínua e a integral (que aqui, aliás, é de Riemann) é um limite, logo a conjugação comuta com a integral, e, da fórmula de inversão,
	\[
		[\overline{\hat\psi}]^\wedge=[(\overline\psi)^\vee]^\wedge=\overline\psi,
	\]
	donde (*) é \(\int\varphi\overline\psi\) e a afirmação segue.

	Em particular,
	\[
		\|\hat\psi\|_{L^2}^2
		=(\hat\psi,\hat\psi)_{L^2}
		=(\psi,\psi)_{L^2}
		=\|\psi\|_{L^2}^2,
	\]
	dizendo que
	\[
		\mathcal{F}:(\mathcal{S},\|\cdot\|_{L^2})\longrightarrow(\mathcal{S},\|\cdot\|_{L^2})
	\]
	é um operador contínuo unitário.

	Como \(\mathcal{S}\) é denso em \(L^2\) (porque \(\mathcal{C}_c^\infty\hookrightarrow\mathcal{S}\) e \(\mathcal{C}_c^\infty\) é denso em \(L^2\)), resulta que \(\mathcal{F}\) possui uma única extensão contínua a \(L^2\), indiquemos
	\[
		\overline{\mathcal{F}}:L^2\longrightarrow L^2
	\]
	essa extensão.

	Note que (\#) permanece válida para \(\overline{\mathcal{F}}\), porque \((\cdot,\cdot)_{L^2}\) é bilinear e sequencialmente contínuo e \(\overline{\mathcal{F}}\) é contínua (verifique).

	Resta mostrar que
	\[
		\left.\overline{\mathcal{F}}\right|_{L^2}=\widetilde{\mathcal{F}},
	\]
	o que pode ser demonstrado como consequência da densidade de \(\mathcal{C}_c^\infty\) em ambos \(L^2\) e \(\mathcal{D}'\), de uma maneira particular; a saber, a sequência \((\psi_j)_j\) de funções teste que verifica a densidade pode ser escolhida como sendo a mesma, por exemplo
	\[
		\chi_j\cdot(\varphi_{\varepsilon_j}*f),
	\]
	com \(\varphi_\varepsilon\) família regularizante e \(\chi_j\) função de Uryhson associada a um esgotamento de \(\mathbb{R}^n\).

	Com efeito, dado \(f\in L^2\), também \(f\in\mathcal{S}'\), seja \((\psi_j)_j\) uma tal sequência em \(\mathcal{C}_c^\infty\). Temos:
	\[
		\text{(I)}\qquad
		\overline{\mathcal{F}}(f)
		=\lim_{j\to\infty}\overline{\mathcal{F}}(\psi_j)
		=\lim_{j\to\infty}\hat\psi_j
		\quad\text{em }L^2,
	\]
	portanto também em \(\mathcal{S}'\), pois \(L^2\hookrightarrow\mathcal{S}'\); e
	\[
		\text{(II)}\qquad
		\widetilde{\mathcal{F}}(f)
		=\lim_{j\to\infty}\widetilde{\mathcal{F}}(\psi_j)
		=\lim_{j\to\infty}\hat\psi_j
		\quad\text{em }\mathcal{S}'.
	\]
	Pela unicidade do limite fraco-* resulta que, em \(\mathcal{S}'\), as duas distribuições coincidem como funcionais definidos em \(\mathcal{S}\):
	\[
		\overline{\mathcal{F}}(f)=\widetilde{\mathcal{F}}(f),
	\]
	completando a prova, porque toda função de \(L^2\) define um único funcional contínuo em \(\mathcal{S}\) e esse é representado por \(\overline{\mathcal{F}}(f)\) (\(\mathcal{S}\to L^2\) denso \(\Rightarrow L^2=[L^2]'\hookrightarrow\mathcal{S}'\), e cada \(g\in L^2\) é, ela própria, um funcional em \(L^2\)).

	Mais precisamente, o funcional \(\widetilde{\mathcal{F}}(f):\mathcal{S}\to\mathbb{C}\) é representado pela função \(\overline{\mathcal{F}}(f)\in L^2\). \qedsymbol
\end{proof*}

\begin{theorem*}[Plancherel]
	\(\overline{\mathcal{F}}:L^2\longrightarrow L^2\) é um operador unitário, sendo a única extensão de \(\mathcal{F}:\mathcal{S}\to\mathcal{S}\) com essa propriedade.
\end{theorem*}

\subsection{Transformada de Fourier--Laplace}

Começamos mencionando, indicando \(\widetilde{\mathcal{F}}(u)=\hat u\), isto é, mudando \(\overline{\mathcal{F}}\) para \(\widetilde{\mathcal{F}}\).

\begin{prop*}
	O mapa \(\widetilde{\mathcal{F}}:\mathcal{S}'\longrightarrow\mathcal{S}'\) possui as propriedades:
	\[
		\text{(I)}\qquad
		\partial_\xi^\alpha\hat u
		=\bigl[(-2\pi i x)^\alpha u\bigr]^\wedge,
		\qquad \forall u\in\mathcal{S}',\ \alpha\in\mathbb{Z}_+^n;
	\]
	\[
		\text{(II)}\qquad
		(\partial_x^\alpha u)^\wedge
		=(2\pi i\xi)^\alpha\hat u,
		\qquad \forall u\in\mathcal{S}',\ \alpha\in\mathbb{Z}_+^n.
	\]
\end{prop*}
\begin{proof*}
	Por exemplo,
	\begin{align*}
		\langle (\partial^\alpha u)^\wedge,\psi\rangle
		 & =\langle \partial^\alpha u,\hat\psi\rangle
		=(-1)^{|\alpha|}\langle u,\partial^\alpha(\hat\psi)\rangle                    \\
		 & =(-1)^{|\alpha|}\bigl\langle u,[(-2\pi i x)^\alpha\psi]^\wedge\bigr\rangle \\
		 & =\langle \hat u,(2\pi i x)^\alpha\psi\rangle
		=\langle (2\pi i\xi)^\alpha\hat u,\psi\rangle,
	\end{align*}
	estabelecendo (II). Analogamente se prova (I). \qedsymbol
\end{proof*}

\begin{tcolorbox}[
		skin=enhanced,
		title=Observação,
		fonttitle=\bfseries,
		colframe=black,
		colbacktitle=cyan!75!white,
		colback=cyan!15,
		colbacklower=black,
		coltitle=black,
		drop fuzzy shadow
	]
	Como \(\mathcal{S}\to\mathcal{C}^\infty\) densamente, resulta que \(\mathcal{E}'\hookrightarrow\mathcal{S}'\); toda distribuição com suporte compacto é transformável e, a seguir, desejamos caracterizar ``essas'' transformadas, o que nos leva à transformada de Fourier--Laplace, cuja definição depende do teorema a seguir.
\end{tcolorbox}

\begin{theorem*}
	Se \(u\in\mathcal{E}'\), então \(\hat u\) é a função dada por
	\[
		\hat u(\xi)=\langle u(x),e^{-2\pi i x\cdot\xi}\rangle,
		\qquad \xi\in\mathbb{R}^n.
	\]
	Além disso, \(\hat u\) admite uma extensão inteira a todo o \(\mathbb{C}^n\), definida pela mesma igualdade
	\[
		\hat u(\zeta):=\langle u(x),e^{-2\pi i x\cdot\zeta}\rangle,
		\qquad \zeta=(\zeta_1,\ldots,\zeta_n)\in\mathbb{C}^n,
	\]
	(vetor complexo), a qual é denominada a \underline{Transformada de Fourier--Laplace de \(u\)}.
\end{theorem*}
\begin{proof*}
	Com efeito, fixemos a família regularizante padrão \((\varphi_\varepsilon)_{\varepsilon>0}\) e \(u\in\mathcal{E}'\) (com \(\varphi_\varepsilon\) par, \(\forall \varepsilon>0\)) e observemos que \(\hat u\in\mathcal{S}'\hookrightarrow\mathcal{D}'\); logo,
	\[
		\hat u*\varphi_\varepsilon\longrightarrow\hat u
		\qquad\text{em }\mathcal{D}',\quad \varepsilon\to0^+.
	\]
	Com isso, é suficiente mostrarmos que \((\hat u*\varphi_\varepsilon)_{\varepsilon>0}\) converge, para a função do enunciado do teorema, também em \(\mathcal{D}'\).

	Para isso, primeiramente, dado \(\xi\in\mathbb{R}^n\), escrevemos
	\begin{align*}
		\hat u*\varphi_\varepsilon(\xi)
		 & =\bigl\langle \hat u,\tau_\xi(\breve\varphi_\varepsilon)\bigr\rangle
		=\bigl\langle \hat u(x),\tau_\xi(\varphi_\varepsilon)(x)\bigr\rangle                 \\
		 & =\bigl\langle u(y),[\tau_\xi\varphi_\varepsilon]^\wedge(y)\bigr\rangle            \\
		 & =\bigl\langle u(y),e^{-2\pi i\xi\cdot y}(\hat\varphi_\varepsilon)(y)\bigr\rangle.
	\end{align*}

	\begin{tcolorbox}[
			skin=enhanced,
			title=Notação,
			fonttitle=\bfseries,
			colframe=black,
			colbacktitle=cyan!75!white,
			colback=cyan!15,
			colbacklower=black,
			coltitle=black,
			drop fuzzy shadow
		]
		Escreveremos \(\breve f(x)=f(-x)\) para não confundirmos reflexão com a transformada inversa.
	\end{tcolorbox}

	Agora, como sabemos,
	\begin{align*}
		(\hat\varphi_\varepsilon)(y)
		 & =\int_{\mathbb{R}^n}e^{-2\pi i y\cdot z}\varepsilon^{-n}\varphi(z/\varepsilon)\,\mathrm{d}z \\
		 & =\int_{\mathbb{R}^n}e^{-2\pi i y\cdot(\varepsilon z)}\varphi(z)\,\mathrm{d}z
		=\hat\varphi(\varepsilon y),
	\end{align*}
	por mudança de variáveis \(z=\varepsilon w\).

	Quando \(|y|<R\), \(R>0\) qualquer, temos, qualquer que seja \(\xi\in\mathbb{R}^n\),
	\[
		\left|e^{-2\pi i\xi\cdot y}-e^{-2\pi i\xi\cdot y}\hat\varphi(\varepsilon y)\right|
		=|1-\hat\varphi(\varepsilon y)|.
	\]
	Sendo \(\hat\varphi\in\mathcal{S}\), dado \(\mu>0\), existe \(\alpha>0\) tal que
	\[
		|z|<\alpha\Longrightarrow |\hat\varphi(z)-\hat\varphi(0)|<\mu.
	\]
	Se \(0<\varepsilon<\alpha/R\) e \(|y|<R\), então \(|\varepsilon y|<\varepsilon R<\alpha\), logo
	\[
		|\hat\varphi(\varepsilon y)-\hat\varphi(0)|<\mu,
	\]
	mas \(\hat\varphi(0)=\int\varphi(z)\,\mathrm{d}z=1\), donde
	\[
		\sup_{\xi\in\mathbb{R}^n}
		\left|e^{-2\pi i\xi\cdot y}-e^{-2\pi i\xi\cdot y}\hat\varphi(\varepsilon y)\right|<\mu,
		\qquad 0<\varepsilon<\alpha/R.
	\]
	Isso mostra que a convergência de \(\hat\varphi(\varepsilon y)e^{-2\pi i y\cdot\xi}\) para \(e^{-2\pi i y\cdot\xi}\), em \(\mathcal{C}_y^\infty(\mathbb{R}^n)\), independe de \(\xi\). Portanto,
	\[
		\hat u*\varphi_\varepsilon(\xi)
		\underset{\varepsilon\to0^+}{\longrightarrow}
		\langle u(x),e^{-2\pi i x\cdot\xi}\rangle
	\]
	em \(\mathcal{C}_y(\mathbb{R}^n)\), em particular em \(\mathcal{D}'\), completando a parte de que \(\hat u\) é, enquanto distribuição, representada pela função contínua
	\[
		f(\xi):=\langle u(x),e^{-2\pi i x\cdot\xi}\rangle.
	\]

	Para demonstrar que \(f\) acima admite uma extensão inteira, calculamos, para \(\zeta\in\mathbb{C}^n\),
	\begin{align*}
		(\#)\qquad e^{-2\pi i x\cdot\zeta}
		 & =\sum_{m=0}^\infty\frac{(-2\pi i x\cdot\zeta)^m}{m!}
		=\sum_{m=0}^\infty\frac{(-2\pi i)^m}{m!}(x\cdot\zeta)^m.
	\end{align*}
	Pelo Teorema Multinomial,
	\[
		(x\cdot\zeta)^m
		=(x_1\zeta_1+\cdots+x_n\zeta_n)^m
		=\sum_{|\alpha|=m}\frac{m!}{\alpha!}x^\alpha\zeta^\alpha.
	\]
	Como
	\[
		e^z=\sum_{m=0}^\infty\frac{z^m}{m!}
	\]
	com convergência em \(\mathcal{C}^\infty(\mathbb{C})\) (vide Conway), se \(K\subset\mathbb{C}^n\) é um compacto qualquer, a convergência em (\#) é uniforme em \(\mathcal{C}^\infty(\mathbb{C})\) para \(\zeta\in K\) (como sabemos da Análise Complexa, para ver isso, é suficiente provar a convergência em (\#) em \(\mathcal{C}(\mathbb{C})\), uniforme para \(\zeta\in K\)).

	\begin{tcolorbox}[
			skin=enhanced,
			title=Observação,
			fonttitle=\bfseries,
			colframe=black,
			colbacktitle=cyan!75!white,
			colback=cyan!15,
			colbacklower=black,
			coltitle=black,
			drop fuzzy shadow
		]
		Na verdade, o argumento a seguir demonstra isso.
	\end{tcolorbox}

	Daí, sendo \(u\in\mathcal{E}'\), temos
	\[
		\hat u(\zeta)=\langle u(x),e^{-2\pi i x\cdot\zeta}\rangle
		=\sum_{m=0}^\infty\frac{(-2\pi i)^m}{m!}
		\left[
			\sum_{|\alpha|=m}\frac{m!}{\alpha!}\langle u(x),x^\alpha\rangle\zeta^\alpha
			\right].\tag{**}
	\]

	Além disso, existem um compacto \(L\subset\mathbb{R}^n\), \(c>0\) e \(r\in\mathbb{N}\) tais que, para todo \(g\in\mathcal{C}^\infty\),
	\[
		|\langle u(x),g\rangle|
		\leq c\sum_{|\beta|\leq r}\sup_{x\in L}|\partial^\beta g(x)|.
	\]
	Em particular, para \(g=x^\alpha\),
	\begin{align*}
		|\langle u(x),x^\alpha\rangle|
		 & \leq c\sum_{|\beta|\leq r}\sup_{x\in L}|\partial^\beta x^\alpha|       \\
		 & =c\sum_{\substack{|\beta|\leq r                                        \\ \beta\leq\alpha}}
		\sup_{x\in L}\left|\frac{\alpha!}{(\alpha-\beta)!}x^{\alpha-\beta}\right| \\
		 & \leq c\sum_{\substack{|\beta|\leq r                                    \\ \beta\leq\alpha}}
		\frac{\alpha!}{(\alpha-\beta)!}M_L^{|\alpha|-|\beta|},
	\end{align*}
	onde
	\[
		M_L:=\sup_{x\in L}|x|\geq1.
	\]
	Consequentemente,
	\[
		|\langle u(x),x^\alpha\rangle|
		\leq cM_L^{|\alpha|}
		\sum_{\substack{|\beta|\leq r\\ \beta\leq\alpha}}
		\frac{\alpha!}{(\alpha-\beta)!}.
	\]

	Com isso, estudamos a convergência absoluta, e uniforme para \(\zeta\in K\), da série em (**):
	\begin{align*}
		 & \sum_{m=0}^\infty\frac{(2\pi)^m}{m!}
		\left[
			\sum_{|\alpha|=m}\frac{m!}{\alpha!}
			|\langle u(x),x^\alpha\rangle|\,|\zeta^\alpha|
		\right]                                                  \\
		 & \leq \sum_{m=0}^\infty\frac{(2\pi)^m}{m!}
		\left[
			\sum_{|\alpha|=m}\frac{m!}{\alpha!}
			\left(cM_L^{|\alpha|}
		\sum_{\substack{|\beta|\leq r                            \\ \beta\leq\alpha}}
			\frac{\alpha!}{(\alpha-\beta)!}\right)|\zeta|^m
		\right]                                                  \\
		 & =\sum_{m=0}^\infty c\frac{(2\pi)^m}{m!}M_L^m|\zeta|^m
		\left[
			\sum_{|\alpha|=m}\frac{m!}{\alpha!}
		\sum_{\substack{|\beta|\leq r                            \\ \beta\leq\alpha}}
			\frac{\alpha!}{(\alpha-\beta)!}
			\right].\tag{***}
	\end{align*}

	\begin{tcolorbox}[
			skin=enhanced,
			title=Observação,
			fonttitle=\bfseries,
			colframe=black,
			colbacktitle=cyan!75!white,
			colback=cyan!15,
			colbacklower=black,
			coltitle=black,
			drop fuzzy shadow
		]
		O que está entre colchetes em (*** ) será menor ou igual que \(2^m a_{(n,r)}n^m.\)
	\end{tcolorbox}

	Agora, fazendo uma estimativa ``grossa'' (que segue por argumentos de contagem), temos:
	\begin{align*}
		\text{(1$^\circ$)}\qquad
		\sum_{|\beta|\leq r}\frac{\alpha!}{(\alpha-\beta)!}
		 & =\sum_{j=0}^r\sum_{|\beta|=j}
		\beta!\frac{\alpha!}{(\alpha-\beta)!\beta!}   \\
		 & \leq \sum_{j=0}^r n\,j!
		\sum_{\substack{|\beta|=j                     \\\beta\leq\alpha}}
		\frac{\alpha!}{(\alpha-\beta)!\beta!}         \\
		 & \leq \sum_{j=0}^r n\,j!
		\sum_{\beta+\gamma=\alpha}
		\frac{\alpha!}{\gamma!\beta!}                 \\
		 & =\sum_{j=0}^r n\,j!\,(2,2,\ldots,2)^\alpha
		=\sum_{j=0}^r n\,j!\,2^{|\alpha|}
		=2^{|\alpha|}a_{(n,r)},
	\end{align*}
	onde \(a_{(n,r)}>0\) só depende de \(n\) e \(r\) (sendo \(r\) dependente apenas de \(u\in\mathcal{E}'\)); e
	\[
		\text{(2$^\circ$)}\qquad
		\sum_{|\alpha|=m}\frac{m!}{\alpha!}
		=(1+1+\cdots+1)^m=n^m,
	\]
	pelo Teorema Multinomial.
	\begin{tcolorbox}[
			skin=enhanced,
			title=Observação,
			fonttitle=\bfseries,
			colframe=black,
			colbacktitle=cyan!75!white,
			colback=cyan!15,
			colbacklower=black,
			coltitle=black,
			drop fuzzy shadow
		]
		O passo (*) vem do Teorema Binomial, generalizado:
		\[
			(x+y)^\alpha=\sum_{\beta+\gamma=\alpha}\frac{\alpha!}{\beta!\gamma!}x^\beta y^\gamma,
		\]
		e usamos \(x=y=(1,1,\ldots,1)\in\mathbb{R}^n\).
	\end{tcolorbox}

	Usando essas estimativas na série (***), vemos que ela é dominada por
	\begin{align*}
		c\sum_{m=0}^\infty\frac{(2\pi M_L|\zeta|)^m}{m!}
		2^m a_{(n,r)}n^m
		 & =c\,a_{(n,r)}\sum_{m=0}^\infty
		\frac{(2\pi M_L\,2n\,|\zeta|)^m}{m!}   \\
		 & =c\,a_{(n,r)}e^{4\pi M_L n|\zeta|},
	\end{align*}
	com convergência uniforme para \(\zeta\) em compactos, completando a demonstração. \qedsymbol
\end{proof*}

\begin{theorem*}[Paley--Wiener]
	Ele traz duas conclusões:
	\begin{itemize}
		\item[(I)] Uma função inteira \(U:\mathbb{C}^n\to\mathbb{C}\) é a transformada de Fourier--Laplace de uma \(\varphi\in\mathcal{C}_c^\infty(\overline{B(0;A)})\) se, e somente se, para cada \(N\in\mathbb{N}\), existe \(C_N>0\) tal que
		      \[
			      |U(\zeta)|\leq C_N(1+|\zeta|)^{-N}e^{2\pi A|\operatorname{Im}(\zeta)|},
			      \qquad \forall \zeta\in\mathbb{C}^n.
		      \]
		\item[(II)] Uma função inteira \(U:\mathbb{C}^n\to\mathbb{C}\) é a transformada de Fourier--Laplace de uma \(u\in\mathcal{E}'(\overline{B(0;A)})\) se, e somente se, existem \(C>0\) e \(N\in\mathbb{N}\cup\{0\}\) tais que
		      \[
			      |U(\zeta)|\leq C(1+|\zeta|)^Ne^{2\pi A|\operatorname{Im}(\zeta)|},
			      \qquad \forall \zeta\in\mathbb{C}^n.
		      \]
	\end{itemize}
\end{theorem*}
\begin{proof*}
	Provaremos somente a parte \((\Rightarrow)\) de (I). Com efeito, seja \(\varphi\in\mathcal{C}_c^\infty(\overline{B(0;A)})\), \(A>0\), e consideremos
	\[
		\hat\varphi(\zeta)
		=\int_{\overline{B(0;A)}}e^{-2\pi i x\cdot\zeta}\varphi(x)\,\mathrm{d}x,
		\qquad \zeta\in\mathbb{C}^n,
	\]
	sua transformada de Fourier--Laplace.

	Por um lado, é fácil ver que, para todo \(\alpha\in\mathbb{Z}_+^n\),
	\[
		(\partial_x^\alpha\varphi)^\wedge(\zeta)
		=(2\pi i\zeta)^\alpha\hat\varphi(\zeta),
	\]
	consequentemente, dados \(N\in\mathbb{Z}_+\) e \(j=1,2,\ldots,n\),
	\[
		(\partial_j^N\varphi)^\wedge(\zeta)
		=(2\pi i\zeta_j)^N\hat\varphi(\zeta),
		\qquad \zeta\in\mathbb{C}^n.
	\]

	Por outro,
	\[
		(\partial_j^N\varphi)^\wedge(\zeta)
		=\int_{\overline{B(0;A)}}e^{-2\pi i x\cdot\zeta}(\partial_j^N\varphi)(x)\,\mathrm{d}x,
	\]
	e, portanto,
	\[
		\left|(\partial_j^N\varphi)^\wedge(\zeta)\right|
		\leq\int_{\overline{B(0;A)}}
		\|\partial_j^N\varphi\|_\infty
		e^{\operatorname{Re}[-2\pi i x\cdot\zeta]}
		\,\mathrm{d}x.\tag{+}
	\]

	Escrevendo
	\[
		\zeta=(\xi_1+i\eta_1,\xi_2+i\eta_2,\ldots,\xi_n+i\eta_n),
	\]
	percebemos que
	\begin{align*}
		i x\cdot\zeta
		 & =i(x_1\xi_1+i x_1\eta_1+x_2\xi_2+i x_2\eta_2+\cdots+x_n\xi_n+i x_n\eta_n) \\
		 & =-(x_1\eta_1+x_2\eta_2+\cdots+x_n\eta_n)
		+i(x_1\xi_1+\cdots+x_n\xi_n),
	\end{align*}
	e, assim,
	\[
		\operatorname{Re}[-2\pi i x\cdot\zeta]=x\cdot\eta,
		\qquad \eta=(\eta_1,\ldots,\eta_n).
	\]
	Donde
	\[
		(+)\leq \|\partial_j^N\varphi\|_\infty
		m(\overline{B(0;A)})e^{2\pi A|\operatorname{Im}(\zeta)|},
	\]
	pois \(|x\cdot\eta|\leq|x|\,|\eta|\leq A|\eta|\), por Cauchy--Schwarz.

	Juntando tudo, para todo \(N\in\mathbb{N}\cup\{0\}\) e \(j=1,2,\ldots,n\),
	\[
		(2\pi)^N|\zeta_j|^N|\hat\varphi(\zeta)|
		=|(\partial_j^N\varphi)^\wedge(\zeta)|
		\leq m(\overline{B(0;A)})\|\partial_j^N\varphi\|_\infty
		e^{2\pi A|\operatorname{Im}(\zeta)|},
	\]
	e, portanto,
	\[
		(2\pi)^N\left(\sum_{j=1}^n|\zeta_j|^N\right)|\hat\varphi(\zeta)|
		\leq\left(\sum_{j=1}^n\|\partial_j^N\varphi\|_\infty\right)
		m(\overline{B(0;A)})e^{2\pi A|\eta|}.
	\]
	Como a função \(t\mapsto t^N\) é convexa e as normas da soma e euclidiana são equivalentes em \(\mathbb{C}^n\), existe \(C_1>0\), que depende só de \(A\) e \(\varphi\), tal que
	\[
		(2\pi)^N|\zeta|^N|\hat\varphi(\zeta)|
		\leq C_1e^{2\pi A|\eta|},
	\]
	e, consequentemente,
	\[
		(2\pi)^N(1+|\zeta|^N)|\hat\varphi(\zeta)|
		\leq C_2e^{2\pi A|\eta|},
	\]
	completando a prova.

	Vamos incluir aqui uma prova da recíproca como cortesia. Com efeito, seja \(U:\mathbb{C}^n\to\mathbb{C}\) inteira tal que
	\[
		|U(\zeta)|\leq C_N(1+|\zeta|)^{-N}e^{2\pi A|\operatorname{Im}(\zeta)|}
	\]
	e consideramos \(U|_{\mathbb{R}^n}\), a qual, pelas estimativas, satisfaz
	\[
		|U(\xi)|\leq C_N(1+|\xi|)^{-N},
		\qquad \forall \xi\in\mathbb{R}^n,
	\]
	donde
	\[
		(1+|\xi|)^N U|_{\mathbb{R}^n}\in L^1,
		\qquad \forall N\in\mathbb{N}.
	\]
	Com isso, podemos definir
	\[
		\varphi(x):=[U|_{\mathbb{R}^n}]^\vee(x)
		=\int_{\mathbb{R}^n}e^{2\pi i x\cdot\xi}U(\xi)\,\mathrm{d}\xi,
		\qquad x\in\mathbb{R}^n.
	\]
	Pelas observações acima, podemos derivar \(\varphi\) sob o sinal da integral indefinidamente, mostrando que \(\varphi\in\mathcal{C}^\infty\).

	Por outro lado, dado \(x\in\mathbb{R}^n\), como \(U\) é inteira e as ``curvas''
	\[
		\gamma(\xi)=\xi
		\qquad\text{e}\qquad
		\lambda(\xi)=\xi+i\eta,
		\qquad \xi\in\mathbb{R}^n,
	\]
	são homotópicas, podemos escrever
	\[
		\varphi(x)=\int_{\mathbb{R}^n}e^{2\pi i x\cdot(\xi+i\eta)}U(\xi+i\eta)\,\mathrm{d}\xi,
	\]
	e, portanto,
	\begin{align*}
		|\varphi(x)|
		 & \leq\int_{\mathbb{R}^n}e^{-2\pi x\cdot\eta}
		C_N(1+|\xi|)^{-N}e^{2\pi A|\eta|}\,\mathrm{d}\xi    \\
		 & \leq C_N\int_{\mathbb{R}^n}(1+|\xi|)^{-N}
		e^{-2\pi x\cdot\eta}e^{2\pi A|\eta|}\,\mathrm{d}\xi \\
		 & =C_N\int_{\mathbb{R}^n}(1+|\xi|)^{-N}
		e^{2\pi(A|\eta|-x\cdot\eta)}\,\mathrm{d}\xi,
	\end{align*}
	qualquer que seja \(\eta\in\mathbb{R}^n\).

	Em particular, escolhendo \(\eta=tx\), \(t>0\), temos
	\[
		x\cdot\eta=t|x|^2
		\qquad\text{e}\qquad
		|\eta|=t|x|,
	\]
	e então
	\[
		e^{2\pi(A|\eta|-x\cdot\eta)}
		=e^{2\pi(At|x|-t|x|^2)}
		=e^{2\pi t|x|(A-|x|)}
		\underset{t\to\infty}{\longrightarrow}0,
	\]
	porque \(|x|>A\). \qedsymbol
\end{proof*}

\end{document}
